\documentclass{article}

% -----------------------------
% Packages
% -----------------------------
\usepackage{fancyhdr}
\usepackage{extramarks}
\usepackage{amsmath}
\usepackage{amsthm}
\usepackage{amsfonts}
\usepackage{tikz}
\usepackage[plain]{algorithm}
\usepackage{algpseudocode}
\usepackage{amssymb}

\usetikzlibrary{automata,positioning}

% -----------------------------
% Basic Document Settings
% -----------------------------
\topmargin=-0.45in
\evensidemargin=0in
\oddsidemargin=0in
\textwidth=6.5in
\textheight=9.0in
\headsep=0.25in

\linespread{1.1}

\pagestyle{fancy}
\lhead{\hmwkAuthorName}
\chead{\hmwkClass: \hmwkTitle}
\rhead{}
\lfoot{\lastxmark}
\cfoot{\thepage}

\renewcommand\headrulewidth{0.4pt}
\renewcommand\footrulewidth{0.4pt}

\setlength\parindent{0pt}

% -----------------------------
% Problem Section Helpers
% -----------------------------
\newcommand{\enterProblemHeader}[1]{
  \nobreak\extramarks{}{Problem #1 continued on next page\ldots}\nobreak{}
  \nobreak\extramarks{Problem #1 (continued)}{Problem #1 continued on next page\ldots}\nobreak{}
}

\newcommand{\exitProblemHeader}[1]{
  \nobreak\extramarks{Problem #1 (continued)}{Problem #1 continued on next page\ldots}\nobreak{}
  \nobreak\extramarks{Problem #1}{}\nobreak{}
}

\setcounter{secnumdepth}{0}
\newcounter{partCounter}
\newcounter{homeworkProblemCounter}
\setcounter{homeworkProblemCounter}{1}
\nobreak\extramarks{Problem \arabic{homeworkProblemCounter}}{}\nobreak{}

% -----------------------------
% Theorem env (optional)
% -----------------------------
\newtheorem*{theorem}{Theorem}

% -----------------------------
% Homework Problem Environment
%   Usage:
%   \begin{homeworkProblem}            % auto-number
%     % Your answer here
%   \end{homeworkProblem}
%
%   \begin{homeworkProblem}[1.2.3]     % label problem "1.2.3" (e.g. textbook numbering)
%     % Your answer here
%   \end{homeworkProblem}
% -----------------------------
\newenvironment{homeworkProblem}[1][]{
  \def\hwProblemLabel{#1}
  \ifx\hwProblemLabel\empty
    \edef\hwProblemLabel{\arabic{homeworkProblemCounter}}
    \stepcounter{homeworkProblemCounter}
  \fi
  \section{Problem \hwProblemLabel}
  \setcounter{partCounter}{1}
  \enterProblemHeader{\hwProblemLabel}
}{
  \exitProblemHeader{\hwProblemLabel}
}

% -----------------------------
% Homework Details (EDIT THESE)
% -----------------------------
\newcommand{\hmwkTitle}{Homework \#1}
\newcommand{\hmwkClass}{Math 3320}
\newcommand{\hmwkAuthorName}{\textbf{Joseph Quinn}}

% -----------------------------
% Title
% -----------------------------
\title{
  \vspace{2in}
  \textmd{\textbf{\hmwkClass:\ \hmwkTitle}}\\
  \vspace{3in}
}
\author{\hmwkAuthorName}
\date{}

% -----------------------------
% Convenience Macros (optional)
% -----------------------------
\renewcommand{\part}[1]{\textbf{\large Part \Alph{partCounter}}\stepcounter{partCounter}\\}

% Algorithms
\newcommand{\alg}[1]{\textsc{\bfseries \footnotesize #1}}

% Calculus
\newcommand{\deriv}[1]{\frac{\mathrm{d}}{\mathrm{d}x} (#1)}
\newcommand{\pderiv}[2]{\frac{\partial}{\partial #1} (#2)}
\newcommand{\dx}{\mathrm{d}x}

% Statistics
\newcommand{\E}{\mathrm{E}}
\newcommand{\Var}{\mathrm{Var}}
\newcommand{\Cov}{\mathrm{Cov}}
\newcommand{\Bias}{\mathrm{Bias}}

% Proof step (for aligned derivations)
\newcommand{\step}[2]{& #1 & & \text{#2} \\}

% Solution header (optional)
\newcommand{\solution}{\textbf{\large Solution}}

% -----------------------------
% Document
% -----------------------------
\begin{document}

\maketitle
\pagebreak

\begin{homeworkProblem}[1.2.2]
  total number of words of length \( n \) =\( \; 2^n \) 
\end{homeworkProblem}

\begin{homeworkProblem}[1.2.3]
  Make it a counting probelm, our sample is the indexes for the word digits, and what we are choosing is the 
  positions that are 1s. So how many different ways can we select a size 2 combinations, a size 4 combination, and a size 6 combination.
  And if we are assuming 0 is even then aso size 0 combinations.
  Then just add them together

  \[
    \begin{pmatrix}
      6 \\ 0
    \end{pmatrix}
    +
    \begin{pmatrix}
      6 \\ 2
    \end{pmatrix}
    + 
    \begin{pmatrix}
      6 \\ 4
    \end{pmatrix}
    +
    \begin{pmatrix}
      6 \\ 6
    \end{pmatrix}
    = 32
  \] 

\end{homeworkProblem}
\begin{homeworkProblem}[1.2.5]
Simply flip the bits of our received codeword making our reliability the complement. 
\end{homeworkProblem}

\begin{homeworkProblem}[1.2.6]
  This is a very bad channel and means every time we send a bit it essentially is just reduced to static and means no error correction is even possible.
\end{homeworkProblem}

\begin{homeworkProblem}[1.3.5]
  (a)
  our code \( C  = \{0000,0011,0101,0110,1001,1010,1100,1111\}\) if we received \( 1101 \) we can detect an error because this word is not in our code.
  \\
  (b)
  The codewords that are the least number of bit swaps away are {1111, 0101, 1001, 1100}
  \\
  (c)
  No because every non codeword of length 4 has an odd number of 1s so by flipping any of the digits it would then have an even number of 1s so no matter
  what you pick it will be 1 flip away from 4 codewords.

\end{homeworkProblem}

\begin{homeworkProblem}[1.3.6]
our code \( C  = \{000000000, 001001001, 010010010, 011011011, 100100100, 101101101, 110110110, 111111111\}\) 
(a): closest to 001001001 \\
(b): closest to 011011011\\
(c): closest to 101101101 \\
(d): closest to 000000000

\end{homeworkProblem}

\begin{homeworkProblem}[1.3.7]
  In class we talked about how parity checks can always detect a single error, so for a parity check on a word of length 3 results in a codeword of length 4,
  this is exactly what we just did and results in a max of 8 codewords.

  This also generalizes to the size of the code for length \( n-1 \), or \( 2^{n-1} \)
\end{homeworkProblem}
\end{document}
